% !TEX program = platex
% 数学ノート集・共通テンプレート v2（積分論2026年度第1部の版面を参考）
% A5判。章内で定義・定理・例・注意等の番号を共通化。
% コンパイル： platex -interaction=nonstopmode math_note_template_v2.tex (2回)
%              dvipdfmx math_note_template_v2.dvi
\documentclass[a5paper,10pt,onecolumn]{jsarticle}
\usepackage[dvipdfmx,top=22mm,bottom=20mm,inner=19mm,outer=18mm,headheight=13pt,headsep=6mm]{geometry}
\usepackage{amsmath,amssymb,amsthm,mathtools}
\usepackage{enumitem,fancyhdr,titlesec}
\usepackage[dvipdfmx,bookmarks=true,colorlinks=false]{hyperref}
\usepackage{pxjahyper}

% ----- 差し替える文書情報 -----
\newcommand{\NoteTitle}{Hilbert–SerreとHilbert–Samuel多項式}
\newcommand{\NoteSeries}{数学ノート集}

% ----- 文字・組版 -----
% 参考PDFの質素な明朝＋Computer Modernの印象を引き継ぎながら，
% A5判で窮屈にならないよう本文の行間を少し確保する．
\linespread{1.13}
\setlength{\parindent}{1zw}
\setlength{\parskip}{0pt}
\setlength{\abovedisplayskip}{6pt plus 2pt minus 2pt}
\setlength{\belowdisplayskip}{6pt plus 2pt minus 2pt}
\setlength{\abovedisplayshortskip}{3pt plus 2pt minus 1pt}
\setlength{\belowdisplayshortskip}{4pt plus 2pt minus 1pt}
\setlist[enumerate]{topsep=3pt,itemsep=1pt,parsep=0pt,leftmargin=2.4zw}
\titleformat{\section}{\normalfont\normalsize\bfseries}{\thesection}{0.8zw}{}
\titlespacing*{\section}{0pt}{2.0ex plus .4ex}{.9ex}
\titleformat{\subsection}{\normalfont\normalsize\bfseries}{\thesubsection}{0.7zw}{}
\titlespacing*{\subsection}{0pt}{1.5ex plus .3ex}{.6ex}
\raggedbottom

% ----- 柱とページ番号：ページ上部、罫線なし -----
\fancypagestyle{mathnote}{%
  \fancyhf{}%
  \fancyhead[L]{\footnotesize \NoteSeries\ （\NoteTitle）}%
  \fancyhead[R]{\footnotesize \thepage}%
  \renewcommand{\headrulewidth}{0pt}%
  \renewcommand{\footrulewidth}{0pt}%
}
\pagestyle{mathnote}

% ----- 節内通し番号の定義・定理群 -----
\newtheoremstyle{mathnotetheorem}
 {4.5pt}{4.5pt}{\normalfont}{}%
 {\bfseries}{}{0.65zw}%
 {\thmname{#1}\thmnumber{ #2}\thmnote{（#3）}}
\theoremstyle{mathnotetheorem}
\newtheorem{theorem}{定理}[section]
\newtheorem{definition}[theorem]{定義}
\newtheorem{lemma}[theorem]{補題}
\newtheorem{proposition}[theorem]{命題}
\newtheorem{corollary}[theorem]{系}
\newtheorem{example}[theorem]{例}
\newtheorem{remark}[theorem]{注意}
\renewcommand{\proofname}{証明}
\renewcommand{\qedsymbol}{$\square$}
\makeatletter
\renewenvironment{proof}[1][\proofname]{%
  \par\pushQED{\qed}\normalfont\topsep4\p@\@plus2\p@\relax
  \trivlist\item[\hskip\labelsep\bfseries #1]\ignorespaces
}{\popQED\endtrivlist\@endpefalse}
\makeatother

% ----- 使うノートで適宜拡張する記号 -----
\newcommand{\R}{\mathbb{R}}
\newcommand{\norm}[1]{\lVert #1\rVert}
\newcommand{\abs}[1]{\lvert #1\rvert}

\begin{document}
\section*{読了後にできること}
本稿を読むための前提は，Artin環，Hilbert級数，局所環と極大イデアル．
本稿の具体的な到達課題は，Hilbert--Samuel関数を具体的な局所環で計算する．

\section{次数付き加群と長さ}
本稿はNoether局所環のHilbert--Samuel函数が，十分大きい整数上で
多項式になることを証明する．出発点は次数付き加群の
Hilbert--Serre型の有理型表示である．

用いる既習事項は，Noether環上の有限生成加群の部分加群も
有限生成であること，Hilbertの基底定理，有限長加群の
Jordan--H\"older長の加法性である．
これらから先の生成函数計算と多項式性は本稿で証明する．

\begin{definition}[次数付き加群]
$R$を可換Artin環とし，
$S=R[T_1,\ldots,T_r]$に$\deg T_i=1$を入れる．
$S$上の\emph{非負次数付き加群}とは
$N=\bigoplus_{n\ge0}N_n$で，
$S_iN_j\subset N_{i+j}$を満たすものをいう．
各$N_n$の$R$-加群としての長さを$\ell_R(N_n)$とする．
\end{definition}

\begin{lemma}[各次数の有限長性]\label{lem:length}
$N$が$S$上有限生成の非負次数付き加群なら，
各$N_n$は有限長$R$-加群である．
\end{lemma}
\begin{proof}
$N$の有限個の斉次生成元$v_1,\ldots,v_m$を取れる．
実際，任意の有限生成元を斉次成分の有限和に分解し，
それらの全斉次成分を生成元に選べばよい．
$n$を固定すると，$N_n$は各$v_i$に乗じる
次数$n-\deg v_i$の単項式の像で$R$上生成される．
該当する単項式は有限個であるから$N_n$は
有限生成$R$-加群である．$R$はArtin環なので
有限生成$R$-加群は有限長となる．
\end{proof}

\begin{definition}[Hilbert級数]
上の$N$のHilbert級数を形式的冪級数
\[
 H_N(z)=\sum_{n\ge0}\ell_R(N_n)z^n\in\mathbb Z[[z]]
\]
で定める．
\end{definition}

\section{Hilbert--Serre型の有理型表示}
\begin{lemma}[長さの完全列公式]\label{lem:exact}
各次数が有限長の次数付き加群について，斉次写像による
次数0の完全列
\[
 0\longrightarrow A\longrightarrow B\longrightarrow C
 \longrightarrow0
\]
があれば$H_B(z)=H_A(z)+H_C(z)$である．
また$N(-1)_n=N_{n-1}$（$n\ge1$），$N(-1)_0=0$
と置けば$H_{N(-1)}(z)=zH_N(z)$である．
\end{lemma}
\begin{proof}
完全列の各次数$n$の部分は
$0\to A_n\to B_n\to C_n\to0$であり，
長さの加法性から$\ell(B_n)=\ell(A_n)+\ell(C_n)$を得る．
これを全$n$について足し合わせれば級数の等式となる．
次数ずらしの公式は添字の置き換えによる．
\end{proof}

\begin{theorem}[Hilbert--Serreの有理型定理]\label{thm:HS}
$R$を可換Artin環，$S=R[T_1,\ldots,T_r]$とし，
$N$を有限生成の非負次数付き$S$-加群とする．
このときHilbert級数は有理型であり，ある整数係数多項式
$P(z)\in\mathbb Z[z]$を用いて
\[
 H_N(z)=\frac{P(z)}{(1-z)^r}
\]
が形式的冪級数として成立する．$r=0$では$H_N$は多項式である．
\end{theorem}
\begin{proof}
$r$について帰納法を行う．$r=0$なら$S=R$は
次数0の環で，有限個の斉次生成元の次数より大きい$n$では
$N_n=0$となる．従って$H_N$は多項式である．

$r\ge1$とし，$T_r$の乗法による次数0の写像
$N(-1)\xrightarrow{T_r}N$を考える．
$K=\{x\in N:T_rx=0\}$，$Q=N/T_rN$と置くと
\[
 0\longrightarrow K(-1)\longrightarrow N(-1)
 \xrightarrow{T_r}N\longrightarrow Q\longrightarrow0
\]
は完全である．
$S$はNoether環なので$K$と$Q$は有限生成次数付き加群である．
しかも$T_rK=T_rQ=0$だから，
両者は$R[T_1,\ldots,T_{r-1}]$上の有限生成次数付き加群となる．
帰納法の仮定から
\[
 H_K(z)=\frac{P_K(z)}{(1-z)^{r-1}},\qquad
 H_Q(z)=\frac{P_Q(z)}{(1-z)^{r-1}}
\]
なる整数係数多項式が存在する．
補題\ref{lem:exact}と上の完全列から
\[
 (1-z)H_N(z)=H_Q(z)-zH_K(z)
\]
となるため，
\[
 H_N(z)=\frac{P_Q(z)-zP_K(z)}{(1-z)^r}
\]
を得る．これで帰納法が完了する．
\end{proof}

\begin{corollary}[Hilbert函数の最終多項式性]\label{cor:eventual}
定理\ref{thm:HS}の仮定のもとで，
$r\ge1$なら整数$n$についての有理係数多項式
$p_N(n)\in\mathbb Q[n]$が存在し，十分大きい$n$に対して
$\ell_R(N_n)=p_N(n)$が成立する．
その次数は高々$r-1$である（零多項式も許す）．
$r=0$なら$N_n=0$が十分大きい$n$で成立する．
\end{corollary}
\begin{proof}
$P(z)=\sum_{j=0}^m a_jz^j$と書き，
\[
 \frac1{(1-z)^r}
 =\sum_{n\ge0}\binom{n+r-1}{r-1}z^n
\]
という形式的冪級数の二項展開を使う．
$n\ge m$では係数比較により
\[
 \ell_R(N_n)=\sum_{j=0}^m
 a_j\binom{n-j+r-1}{r-1}
\]
であり，右辺は$n$の有理係数多項式で次数高々$r-1$である．
$r=0$は定理\ref{thm:HS}の最初の場合に従う．
\end{proof}

\section{Noether局所環と随伴次数環}
$A$を可換Noether局所環，$\mathfrak m$をその極大イデアルとする．
$M$を有限生成$A$-加群，$I\subset A$を$\mathfrak m$-一次イデアル
（$\sqrt I=\mathfrak m$）とする．

\begin{lemma}[Artin商と有限長性]\label{lem:artin}
ある$N\ge1$について$\mathfrak m^N\subset I$となる．
$A/I$はArtin環であり，任意の有限生成$A/I$-加群は
$A$-加群として有限長である．
\end{lemma}
\begin{proof}
$\mathfrak m=(x_1,\ldots,x_q)$と有限生成できる．
$\sqrt I=\mathfrak m$だから各$j$について
$x_j^{e_j}\in I$となる自然数$e_j$が存在する．
$N=1+\sum_j(e_j-1)$と置くと，次数$N$の任意の
生成単項式$x_1^{a_1}\cdots x_q^{a_q}$では
ある$j$について$a_j\ge e_j$である．よって
$\mathfrak m^N\subset I$を得る．

$A/\mathfrak m^N$は
\[
 A/\mathfrak m^N\supset\mathfrak m/\mathfrak m^N
 \supset\cdots\supset\mathfrak m^{N-1}/\mathfrak m^N
 \supset0
\]
という有限濾過を持ち，各商
$\mathfrak m^j/\mathfrak m^{j+1}$は
剰余体$k=A/\mathfrak m$上有限次元である．
実際，$\mathfrak m^j$は有限生成で，商は$\mathfrak m$に
消されるからである．従って$A/\mathfrak m^N$は
有限長であり，その商$A/I$も有限長である．
有限長環はArtin環である．その上の有限生成加群は
有限個の有限長環の直和の商なので有限長となる．
\end{proof}

\begin{definition}[随伴次数環と加群]
\[
 \operatorname{gr}_I(A)=\bigoplus_{n\ge0}I^n/I^{n+1},
 \qquad
 \operatorname{gr}_I(M)=\bigoplus_{n\ge0}I^nM/I^{n+1}M
\]
と定める．乗法および加群作用は代表元の積で定める．
各次数$n$の部分は$I$に消されるため$A/I$上の加群である．
\end{definition}

\begin{lemma}[随伴次数加群の有限生成性]\label{lem:gr}
$I=(a_1,\ldots,a_r)$とすると，
$\operatorname{gr}_I(M)$は
$(A/I)[T_1,\ldots,T_r]$上の有限生成次数付き加群である．
ただし$T_i$は$a_i\bmod I^2$の作用に対応する．
\end{lemma}
\begin{proof}
$\operatorname{gr}_I(A)$は$A/I$と
$a_1+I^2,\ldots,a_r+I^2$によって次数付き環として生成される．
これは$I^n$が長さ$n$の$a_i$の積の$A$-線形結合で
生成されることから従う．従って次数付き全射
\[
 (A/I)[T_1,\ldots,T_r]\longrightarrow
 \operatorname{gr}_I(A)
\]
がある．$M$の有限生成元$m_1,\ldots,m_q$を選べば，
$\operatorname{gr}_I(M)$はその次数0の像によって
$\operatorname{gr}_I(A)$上生成される．
よって上の多項式環上でも有限生成である．
\end{proof}

\section{Hilbert--Samuel多項式の存在}
\begin{theorem}[Hilbert--Samuel多項式]\label{thm:HSamuel}
$A$をNoether局所環，$M$を有限生成$A$-加群，
$I=(a_1,\ldots,a_r)$を$\mathfrak m$-一次イデアルとする．
このときある有理係数多項式$Q_{I,M}(n)$が存在して，
十分大きい$n$に対して
\[
 \ell_A(M/I^{n+1}M)=Q_{I,M}(n)
\]
となる．その次数は高々$r$である．
\end{theorem}
\begin{proof}
補題\ref{lem:artin}と補題\ref{lem:gr}により，
$G=\operatorname{gr}_I(M)$はArtin環$R=A/I$上の
$r$変数多項式環の有限生成次数付き加群である．
各$G_n=I^nM/I^{n+1}M$は有限長であり，
$A$-加群としての長さと$R$-加群としての長さは一致する．
定理\ref{thm:HS}を適用すると
\[
 \sum_{n\ge0}\ell_A(I^nM/I^{n+1}M)z^n
 =\frac{P(z)}{(1-z)^r}
\]
なる$P(z)\in\mathbb Z[z]$が存在する．

短完全列と長さの加法性を繰り返すことで
\[
 \ell_A(M/I^{n+1}M)
  =\sum_{j=0}^n\ell_A(I^jM/I^{j+1}M)
\]
を得る．したがって形式的冪級数として
\[
 \sum_{n\ge0}\ell_A(M/I^{n+1}M)z^n
 =\frac{1}{1-z}\frac{P(z)}{(1-z)^r}
 =\frac{P(z)}{(1-z)^{r+1}}.
\]
$r\ge1$なら系\ref{cor:eventual}を
$r+1$変数の場合の係数展開として用いれば，
係数は十分大きい$n$で次数高々$r$の多項式となる．
$r=0$なら$I=0$であり，$\sqrt I=\mathfrak m$ゆえ
$A$はArtin局所環で，$\ell_A(M)<\infty$である．
そのとき左辺は定数$\ell_A(M)$であるから結論は成立する．
\end{proof}

\begin{example}[正則局所環の最も基本的な場合]
$A=k[[x_1,\ldots,x_r]]$，$\mathfrak m=(x_1,\ldots,x_r)$，
$M=A$とする．次数が$n$以下の単項式の個数を数えると
\[
 \ell_A(A/\mathfrak m^{n+1})
   =\binom{n+r}{r}.
\]
特にこの場合のHilbert--Samuel多項式の次数は$r$に等しい．
この等式は形式的冪級数の単項式基底から従い，
一般の局所環におけるKrull次元との一致を証明するものではない．
\end{example}

\section{本稿の到達範囲}
\begin{remark}
本稿はHilbert--Samuel函数の最終多項式性と
次数の生成元数による上界を証明した．
さらに強い主張
\[
 \deg Q_{I,M}=\dim_A M
\]
（適切な非零加群と次元の規約の下で成立する）には，
局所環の次元論やHilbert--Samuel重複度の理論が必要である．
この同一性を本稿の結論として暗黙に使用しない．
\end{remark}

\begin{thebibliography}{9}
\bibitem{atiyah} M. F. Atiyah and I. G. Macdonald,
 \emph{Introduction to Commutative Algebra}, Addison-Wesley, 1969.
\bibitem{matsumura} H. Matsumura,
 \emph{Commutative Ring Theory}, Cambridge University Press, 1986.
\end{thebibliography}

\section{到達課題の完全解答}
\noindent\textbf{Hilbert--Samuel関数の完全な計算．}
\begin{proposition}
$R=k[[x,y]]$，$\mathfrak m=(x,y)$ とすると，
全ての $n\ge0$ に対し
\[
\operatorname{length}_R(R/\mathfrak m^{n+1})
=\binom{n+2}{2}
\]
である．
\end{proposition}
\begin{proof}
$\mathfrak m^{n+1}$ は全次数 $n+1$ 以上の
単項式を含むため，商の $k$-基底は
$x^ay^b$（$a,b\ge0$，$a+b\le n$）の剰余類である．
固定された全次数 $j$ の単項式は $j+1$ 個だから
基底数は
$\sum_{j=0}^n(j+1)=(n+1)(n+2)/2$．
この商は剰余体 $k$ を組成因子とする有限長加群であり，
その長さは $k$-次元に等しい．
\end{proof}
\begin{corollary}
$R$ のHilbert--Samuel多項式は
$(n+1)(n+2)/2$ であり，重複度は1である．
\end{corollary}
\begin{proof}
上の式は全ての $n\ge0$ で成立する二次多項式である．
最高次項は $n^2/2$ だから，
次元2の局所環の重複度の定義より $e(R)=1$．
\end{proof}
\noindent\textbf{到達確認．}
$k[[x]]$ で $\operatorname{length}(R/(x^{n+1}))$
を求めよ．
\par\noindent\textbf{解答．}
$1,x,\dots,x^n$ の剰余類が基底なので $n+1$ である．
\end{document}